\documentclass[11pt]{article}
\usepackage{amsmath,amssymb,amsthm,mathtools,hyperref}

\begin{document}

\begin{quote}
This manuscript gives an attempted answer for each of the 330 supplied catalogue
entries. It preserves the original numbering and states the target beside the
attempt, distinguishing a theorem about a specified model from resolution of a
biological research target. The results consist of worked special cases,
identifiability obstructions, quantitative bounds, and explicitly preliminary
reductions. None is claimed to resolve its complete catalogue target. No new
biological experiments, empirical data analyses, or external proof verification
were performed for this manuscript.
\end{quote}

\section*{How to read the attempts}
\addcontentsline{toc}{section}{How to read the attempts}
\paragraph{Scope of the mathematical claims.}
Every proposition applies only under its displayed assumptions. The proof is an
argument for that proposition, not a claim that a biological system satisfies
those assumptions. A known finite generator, validated loss distribution,
identified parameter, or measured response function is an input when the text
says it is known. Computing with such an input does not solve its identification.
The calculations are not claimed to be new discoveries; many are standard
mathematical baselines or direct counterexamples.

\paragraph{Uneven depth and preliminary entries.}
There are 200 saved problem-directed attempts and 130 preliminary conditional reductions. Problems 201--330 are preliminary reductions: a worked model or statistical criterion is provided, but its applicability to the full catalogue formulation has not been established. These entries receive 5\%, not the score of a solved scientific subproblem. 
The later preliminary entries should be developed further before being treated
as research-level attempted solutions. This distinction is also marked locally.

\paragraph{Completion Estimate.}
Percentages are conservative, subjective assessments of coverage of the
\emph{full stated target}, not probabilities of correctness, measured fractions
of a proof, forecasts, or literature-wide progress estimates. In particular,
100\% would require resolving every part of the chosen target under its stated
assumptions, including empirical validation where it is part of the target.
A complete proof for a narrower model does not earn 100\% for the original
question. No entry in this manuscript is assigned 100\%.

\paragraph{Sources and current-status limits.}
Problem statements, definitions, alternative formulations and local source lists
come from the supplied files
\texttt{biology\_open\_problems\_300(2).tex} and
\texttt{biology\_open\_problems\_301\_330(2).tex}.
The organization follows the supplied attempt template
\texttt{O33-23\_v3.tex}: scope, substantive calculation, explicit gap,
completion estimate, and references. Inherited source annotations and links are
catalogue provenance, not evidence that every link or current open status has
been independently rechecked. Selected methodological references appear at the
end. No exhaustive literature or novelty audit is claimed.

\paragraph{Notation and corrections.}
Notation is local to each problem and is restated when a restricted model uses
additional variables. A probability model, metric, loss and observation protocol
must be specified before its numerical quantities are meaningful. One obvious
missing backslash in the reconstruction norm of Problem~77 has been repaired.
No biological mechanism is silently inferred from that typographical repair.

\paragraph{Use limitations.}
Theoretical examples involving disease, therapies, genetic systems, or ecological
interventions are not operational protocols, medical guidance, deployment advice,
or demonstrated safety guarantees. Their assumptions and empirical gaps are part
of the result, not optional qualifications.


\section*{Left–right asymmetry}
\label{prob:129}
\noindent\textbf{Status: restricted result or obstruction; full target not established.}
\subsection{Target and scope}
\paragraph{Catalogue definitions and assumptions.}
Let $q(t)$ be signed fluid flow near the embryonic organiser, $s(t)$ a signed asymmetric signalling variable, and $Y\in\{-1,1\}$ final organ laterality. A measured chiral bias $h$ affects flow; signalling follows a finite stochastic feedback model $ds=f_\theta(s,q,h)dt+\sqrt{2D}\,dW_t$. Define normal sign $Y=1$, noise intensity $D$ and developmental decision time $T$. Mechanical and chemical flow sensing are competing models rather than assumed equivalent.
\paragraph{Original target.}
Predict laterality errors and rescue under changed cilia number, flow direction or signal gain:
\[
Y=\operatorname{sign}s(T),\qquad
\mathbb{P}(Y=-1\mid h,\theta,D)\leq\delta.
\]
Determine the minimal bias and feedback strength consistent with measured reliability, including the full distributions of flow and signalling. Explain how a small physical asymmetry becomes a robust organismal decision.
\paragraph{Catalogue alternative formulations.}
(i) In a ciliated mouse organiser, compare mechanical and transported-chemical sensing using independently manipulated flow and molecule transport. (ii) In a species without an equivalent ciliated organiser, test cell-intrinsic chirality coupled to tissue feedback. These are separate mechanisms whose shared reduced sign dynamics do not imply identical molecular implementation.
\subsection{Attempted answer}
\begin{proposition}
For the linear biased amplifier $ds=(\lambda s+h)dt+\sqrt{2D}\,dW$, $s(0)=0$, $\lambda,D>0$, the sign-error probability at $T$ is $\Phi[-h(e^{\lambda T}-1)/\sqrt{D\lambda(e^{2\lambda T}-1)}]$. Thus error at most $\delta<1/2$ requires $h\geq\Phi^{-1}(1-\delta)\sqrt{D\lambda(e^{2\lambda T}-1)}/(e^{\lambda T}-1)$.
\end{proposition}
\begin{proof}
The integrating-factor solution is Gaussian with mean $h(e^{\lambda T}-1)/\lambda$ and variance $D(e^{2\lambda T}-1)/\lambda$. Standardization gives the tail probability and inversion gives the bias bound.
\end{proof}
\subsection{What remains unresolved}
Linear amplification boosts noise too; the required bias does not vanish at long times in this model. Nonlinear commitment, measured chirality and physical sensing mechanisms remain unvalidated.
\subsection{Completion Estimate}
\textbf{25\%.} The displayed calculation is established under its additional assumptions. The broader theoretical and biological requirements listed above are not established.
\subsection{References}
\noindent\emph{Catalogue sources, retained as provenance:}
\begin{itemize}
\item Saunders and Ingham, Open questions: how to get developmental biology into shape? (2019). \url{https://link.springer.com/article/10.1186/s12915-019-0636-6} (2019 research agenda; supports morphogenesis broadly.)
\item Yoshiba et al., Cilia at the node of mouse embryos sense fluid flow for left-right determination via Pkd2 (2012). \url{https://pubmed.ncbi.nlm.nih.gov/22983710/} (Primary cilia intervention study; indexed abstract checked.)
\end{itemize}




\section{Shared mathematical bounds}\label{app:bounds}
\begin{lemma}[Independent bounded sums]\label{lem:bounded}
Let independent real variables $X_i$ satisfy $a_i\leq X_i\leq b_i$ almost surely,
and let $S=\sum_{i=1}^n(X_i-\mathbb{E} X_i)$. If
$V=\sum_i(b_i-a_i)^2>0$, then for every $t>0$,
\[
 \mathbb{P}(|S|\geq t)\leq 2\exp(-2t^2/V).
\]
In particular, independent losses in $[0,1]$ have mean-error bound
$\mathbb{P}(|\overline X-\mathbb{E}\overline X|\geq\epsilon)
\leq2e^{-2n\epsilon^2}$. The same expression for variables in $[-1,1]$
is $2e^{-n\epsilon^2/2}$. If $V=0$ then $S=0$ almost surely.
\end{lemma}
\begin{proof}
For $X\in[a,b]$ with mean $\mu$, the nonnegative product
$(X-a)(b-X)$ gives
$\operatorname{Var} X\leq(\mu-a)(b-\mu)\leq(b-a)^2/4$.
Let $K(s)=\log\mathbb{E}\exp[s(X-\mathbb{E} X)]$.
Boundedness permits two differentiations. Its second derivative is the variance
under the exponentially tilted law, which has the same support, so
$K''(s)\leq(b-a)^2/4$ for every real $s$.
Because $K(0)=K'(0)=0$, integration yields
$K(s)\leq s^2(b-a)^2/8$.
Independence adds these log moment-generating bounds for $S$.
For $s>0$, Markov's inequality gives
$\mathbb{P}(S\geq t)\leq\exp(-st+s^2V/8)$.
Choosing $s=4t/V$ yields the upper-tail bound $e^{-2t^2/V}$.
Apply it also to $-S$ and add. The stated mean versions follow by substituting
$t=n\epsilon$ and the respective support widths.
\end{proof}

\paragraph{Finite-family extension.}
For $M$ prespecified bounded-loss comparisons, apply the relevant inequality to
each and use a union bound. Independence between comparisons is not required;
independence of the biological units within each comparison is. Repeated cells,
reads or time points from one dependent biological unit cannot automatically be
counted as independent units.

\paragraph{No data-dependent shortcut.}
The bound does not justify choosing a model or validation loss using the same
outcomes and then reporting single-model coverage. A prespecified finite family,
a separate validation sample, or an appropriate complexity or sequential argument
is needed. Marginal predictive coverage likewise should not be read as coverage
for every subgroup or under arbitrary distribution shift.


\section{Methodological references}
The references below support general methods. The detailed subject references
are retained locally with each problem. These entries are not a claim that the
preceding problem-directed calculations appear as new results in these sources.
\begin{thebibliography}{9}
\bibitem{hoeffding}
W. Hoeffding. Probability inequalities for sums of bounded random variables.
\emph{Journal of the American Statistical Association} 58 (1963), 13--30.
\url{https://doi.org/10.1080/01621459.1963.10500830}.
The relevant inequality is proved above rather than left as an external step.

\bibitem{pearl}
J. Pearl. Causal inference in statistics: An overview.
\emph{Statistics Surveys} 3 (2009), 96--146.
\url{https://doi.org/10.1214/09-SS057}.

\bibitem{lei}
J. Lei, M. G'Sell, A. Rinaldo, R. J. Tibshirani, and L. Wasserman.
Distribution-Free Predictive Inference For Regression. Preprint (2016).
\url{https://arxiv.org/abs/1604.04173}.

\bibitem{briat}
C. Briat, A. Gupta, and M. Khammash.
Antithetic Integral Feedback ensures robust perfect adaptation in noisy
biomolecular networks. Preprint (2014).
\url{https://arxiv.org/abs/1410.6064}.
\end{thebibliography}

\end{document}
